\begin{theorem}[Binding and graded agreement of the GBCA specification]
  \label{thm:gbca-spec-binding}
  \lean{PLTS.ABA.GBCA.excluded_card_le_one}\lean{PLTS.ABA.GBCA.retG_value_agree}\lean{PLTS.ABA.GBCA.retG_grade_exclusive}\lean{PLTS.ABA.GBCA.retG_bound_agree}\lean{PLTS.ABA.GBCA.retG_value_eq_bound}\lean{PLTS.ABA.GBCA.retGrade0_excluded_nonempty}\leanok
  \uses{def:aba-gbca-spec}
  Fix an execution of the round-$r$ GBCA specification instance.
  Any two of its returns that hand out a bit hand out the same bit, whatever their grades and whichever processes they answer; the bit a graded outcome hands out is the one $(b,2)$ and $(b,1)$ name, $(\bot,0)$ naming none.\lean{PLTS.ABA.GBCA.outValue}\leanok{} No two of its returns are a grade-2 return and a grade-0 return, which is the second clause of graded agreement.\lean{PLTS.ABA.GBCA.retG_grade_exclusive}\leanok{} Read off the trace, a grade-2 return in a positive-mass trace of the instance excludes every grade-0 return from it.\lean{PLTS.ABA.GBCA.specInst_grade_agree}\leanok{} Any two of its returns announce the same bit.\lean{PLTS.ABA.GBCA.retG_bound_agree}\leanok{} A return that hands out $v$ announces $v$.\lean{PLTS.ABA.GBCA.retG_value_eq_bound}\leanok{} The exclusion set holds at most one bit at every state of the execution, so its reachable shape is $excluded = \emptyset$ or $excluded = \{b\}$ (D19).\lean{PLTS.ABA.GBCA.excluded_card_le_one}\leanok{} A grade-0 return announcing $bnd$ moreover leaves $excluded$ nonempty in every later state, $1-bnd$ being the member, so $1-bnd$ is a bit no extension of the run hands out at grade $\ge 1$ and $bnd$ is the Graded Binding witness.\lean{PLTS.ABA.GBCA.retGrade0_excluded_nonempty}\leanok{}
\end{theorem}

\begin{proof}\leanok
  \uses{def:aba-gbca-spec}
  The exclusion set never shrinks, its single writer inserting and corruption leaving it alone, so the inclusion holds rule by rule and lifts to whole executions.\lean{PLTS.ABA.GBCA.Step.excluded_mono}\leanok{} The cardinality clause is an induction along the execution, the field starting empty and its single writer firing from $\emptyset$ alone.\lean{PLTS.ABA.GBCA.Step.excluded_card_le_one}\leanok{} Inverting the two value-bearing return rows yields the guard pair $v \notin excluded$, $1-v \in excluded$ at the state each fires from, and inverting any return row yields $1-bnd \in excluded$ at its own state.\lean{PLTS.ABA.GBCA.retG_value_guards}\lean{PLTS.ABA.GBCA.retG_bound_guard}\leanok{} For the values: monotonicity carries $1-v_1 \in excluded$ from the earlier return to the later, which demands $v_2 \notin excluded$, so $v_2 \neq 1-v_1$ and on a two-element domain $v_2 = v_1$.
  For the announcements: monotonicity carries $1-bnd_1 \in excluded$ to the state of the later return, which also excludes $1-bnd_2$, and no state of an execution excludes two bits, so $bnd_1 = bnd_2$.
  For the two together: at the state of one value-bearing return the value guard excludes $1-v$ and the bound guard excludes $1-bnd$, and the same cardinality bound gives $v = bnd$.
  The grade-0 return clause is the same monotonicity applied to the witness its own guard produces.
  For the grades: the field $grade$ is written by the grade-2 return alone to the grade-2 side and by the grade-0 return alone to the grade-0 side, each firing only from a state whose grade is unset or already on its own side, and no other rule touches it, so a written grade persists along the execution and the other side's return finds its guard refuted.\lean{PLTS.ABA.GBCA.Step.grade_mono}\lean{PLTS.ABA.GBCA.grade_stable}\leanok{}
  No invariant, no quorum arithmetic and no reachability hypothesis beyond membership in one execution enters (D19).
\end{proof}
